% ===== PRÉAMBULE CANONIQUE — à coller VERBATIM en tête de CHAQUE .tex =====
\documentclass[11pt,a4paper]{article}
\usepackage[utf8]{inputenc}\usepackage[T1]{fontenc}\usepackage{lmodern}
\usepackage[french]{babel}
\usepackage{amsmath,amssymb,mathtools}\usepackage{icomma}
\usepackage[version=4]{mhchem}          % \ce{...} : équations & formules
\usepackage{siunitx}\sisetup{locale=FR} % \qty{}{}, \si{}, \ang{}
\usepackage{chemfig}                    % \chemfig{...} : structures développées
\usepackage{xcolor}\usepackage{colortbl}
\usepackage{tikz}\usepackage{pgfplots}\pgfplotsset{compat=1.18}
\usepackage[most]{tcolorbox}\usepackage{enumitem}\usepackage{array,booktabs}
\usepackage[a4paper,margin=2.2cm]{geometry}
\usetikzlibrary{arrows.meta,calc,angles,quotes,positioning,patterns,backgrounds,shapes.geometric}
\definecolor{primary}{RGB}{222,49,99}\definecolor{primarydark}{RGB}{150,22,55}
\definecolor{defcol}{RGB}{0,114,178}\definecolor{methcol}{RGB}{52,92,148}
\definecolor{excol}{RGB}{0,135,90}\definecolor{attncol}{RGB}{200,40,55}
\tcbset{boxbase/.style={enhanced,breakable,boxrule=0.9pt,arc=2.5pt,
  left=3mm,right=3mm,top=2.3mm,bottom=2.3mm,fonttitle=\bfseries,coltitle=white}}
% --- boîtes TUTORIEL (noms EXACTS, ne pas renommer) ---
\newtcolorbox{cle}[1][]{boxbase,colback=defcol!6,colframe=defcol,
  title={\textsc{À retenir}\ifx\relax#1\relax\relax\else\textmd{ : \itshape #1}\fi}}
\newtcolorbox{methode}[1][]{boxbase,colback=methcol!7,colframe=methcol,sharp corners,
  title={\textsc{Méthode}\ifx\relax#1\relax\relax\else\textmd{ --- \itshape #1}\fi}}
\newtcolorbox{exemplebox}[1][]{boxbase,colback=excol!5,colframe=excol,
  title={\textsc{Exemple}\ifx\relax#1\relax\relax\else\textmd{ : \itshape #1}\fi}}
\newtcolorbox{piege}[1][]{boxbase,colback=attncol!6,colframe=attncol,
  title={\textsc{Piège}\ifx\relax#1\relax\relax\else\textmd{ : \itshape #1}\fi}}
\newtcolorbox{remarque}[1][]{boxbase,colback=black!4,colframe=black!45,
  title={\textsc{Remarque}\ifx\relax#1\relax\relax\else\textmd{ : \itshape #1}\fi}}
% --- boîtes EXERCICE / CORRIGÉ (noms EXACTS, auto-comptées) ---
\newtcolorbox[auto counter]{exo}[1][]{boxbase,colback=primary!4,colframe=primary,
  title={\textsc{Exercice}~\thetcbcounter\ifx\relax#1\relax\relax\else\textmd{ --- \itshape #1}\fi}}
\newtcolorbox[auto counter]{corrige}[1][]{boxbase,colback=excol!4,colframe=excol,
  title={\textsc{Corrigé}~\thetcbcounter\ifx\relax#1\relax\relax\else\textmd{ --- \itshape #1}\fi}}
\newcommand{\R}{\mathbb{R}}\newcommand{\N}{\mathbb{N}}\newcommand{\Z}{\mathbb{Z}}\newcommand{\ds}{\displaystyle}
% =========================================================================
\begin{document}
\begin{corrige}[signe de $\Delta S$ d'après les gaz]
On compare le nombre de moles de gaz réactifs $\to$ produits.
\begin{enumerate}[label=\alph*)]
  \item $4$ mol de gaz $\to 2$ mol : le désordre diminue, $\Delta S<0$.
  \item $0$ mol de gaz $\to 1$ mol : on crée du gaz, $\Delta S>0$.
  \item $3$ mol de gaz $\to 2$ mol : le désordre diminue, $\Delta S<0$.
  \item $2$ mol de gaz $\to 4$ mol : le désordre augmente, $\Delta S>0$.
\end{enumerate}
\end{corrige}

\begin{corrige}[changements d'état]
\begin{enumerate}[label=\alph*)]
  \item solide $\to$ liquide : plus de désordre, $\Delta S>0$.
  \item liquide $\to$ gaz : forte augmentation du désordre, $\Delta S>0$.
  \item gaz $\to$ liquide : le désordre diminue, $\Delta S<0$.
  \item solide $\to$ gaz : forte augmentation du désordre, $\Delta S>0$.
\end{enumerate}
\end{corrige}

\begin{corrige}[dissolution et précipitation]
\begin{enumerate}[label=\alph*)]
  \item un solide ordonné se disperse dans l'eau : $\Delta S>0$ (dissolution).
  \item des ions dispersés forment un solide ordonné : $\Delta S<0$ (précipitation).
\end{enumerate}
\end{corrige}

\begin{corrige}[ordre des états et 3\textsuperscript{e} principe]
\begin{enumerate}[label=\alph*)]
  \item par entropie croissante : \; solide $<$ liquide $<$ gaz.
  \item d'après le 3\textsuperscript{e} principe, $S = \qty{0}{\joule\per\mol\per\kelvin}$ à $T=\qty{0}{\kelvin}$.
\end{enumerate}
\end{corrige}

\begin{corrige}[le 2\textsuperscript{e} principe]
\begin{enumerate}[label=\alph*)]
  \item un processus est spontané si $\Delta S_{\text{univers}} > 0$ (le désordre total augmente).
  \item il fait \textbf{croître} le désordre de l'univers.
  \item non : le 2\textsuperscript{e} principe donne le \emph{sens} d'évolution, pas la \emph{vitesse} (spontané $\ne$ rapide).
\end{enumerate}
\end{corrige}

\end{document}
